\documentclass[aspectratio=169,10pt,spanish]{beamer}

\input{../../../_preambulo_beamer.tex}
\input{../../../_comandos_beamer.tex}

\selectlanguage{spanish}

\title{MA1001B - Modelación Estadística para la Toma de Decisiones}
\subtitle{Lección 34: Prueba z para una media con varianza conocida}
\author{}
\institute{Tecnol\'ogico de Monterrey}
\date{}

\begin{document}

\begin{frame}[plain]
  \vspace{1.2cm}
  {\Large\bfseries MA1001B - Modelación Estadística para la Toma de Decisiones\par}
  \vspace{0.7cm}
  {\large Lección 34: Prueba z para una media con varianza conocida\par}
  \vspace{0.7cm}
  {\color{TecRojo}\rule{\textwidth}{1.2pt}}
  \vspace{0.8cm}

  Facultad MA1001B\\[0.3cm]
  Periodo\\[0.3cm]
  Tecnol\'ogico de Monterrey
\end{frame}

\begin{frame}{La pregunta de la prueba con sigma conocida}
  \begin{alertblock}{Pregunta guía}
    ¿Cómo se lleva a cabo una prueba unilateral de una media desconocida
    cuando $\sigma$ se trata como conocida?
  \end{alertblock}

  \vspace{0.6em}
  Al final de esta lección, usted puede:
  \begin{itemize}
    \item calcular $z=(\bar{x}-\mu_{0})/(\sigma/\sqrt{n})$;
    \item obtener el valor $p$ de cola superior de la normal estándar;
    \item decidir con $\alpha=0.05$ comparando $p$ con $\alpha$;
    \item explicar por qué una $\sigma$ conocida rara vez es verdadera en
      datos de operaciones.
  \end{itemize}

  \vfill
  \scriptsize Todos los tiempos de entrega usados hoy son sintéticos. No se usan datos reales de empresa.
\end{frame}

\begin{frame}{Una prueba sintética de tiempo de entrega}
  \small
  \begin{center}
    \begin{tabular}{lr}
      \toprule
      \textbf{Cantidad} & \textbf{Valor} \\
      \midrule
      $H_{0}$ & $\mu=80$ minutos \\
      $H_{1}$ & $\mu>80$ minutos \\
      $n$, $\sigma$ conocida & 40, 6 minutos \\
      $\bar{x}$ muestral (semilla 42) & 82.229590 minutos \\
      \bottomrule
    \end{tabular}
  \end{center}

  \vspace{0.5em}
  \begin{exampleblock}{Error estándar}
    \[
      \mathrm{EE}=\frac{6}{\sqrt{40}}=0.948683.
    \]
    El analista trata $\sigma$ como conocida a partir de la capacidad de
    proceso de largo plazo.
  \end{exampleblock}
\end{frame}

\begin{frame}[fragile]{Paso 1: El estadístico z}
  \begin{columns}[T]
    \begin{column}{0.42\textwidth}
      \small
      \begin{lstlisting}
se = sigma / np.sqrt(n)
z = (xbar - 80) / se
      \end{lstlisting}

      \begin{block}{Salida verificada}
        $\bar{x}=82.229590$\\
        $\mathrm{EE}=0.948683$\\
        $z=2.350194$
      \end{block}
    \end{column}
    \begin{column}{0.58\textwidth}
      \centering
      \includegraphics[width=\textwidth]{../figuras/prueba_z_media_01_estadistico_z.png}
      \vspace{0.2em}
      \scriptsize Histograma de entregas del Paso 1
    \end{column}
  \end{columns}
\end{frame}

\begin{frame}{Valor $p$ de cola superior y decisión}
  \small
  \[
    p=P(Z\ge 2.350194)=\texttt{norm.sf}(2.350194)=0.009382.
  \]
  Como $p<\alpha=0.05$, se rechaza $H_{0}$: la muestra sintética es
  inconsistente con una media de 80 minutos cuando $\sigma=6$ se trata como
  conocida.

  \vspace{0.4em}
  El valor $p$ no es la probabilidad de que $H_{0}$ sea verdadera. Es la
  probabilidad de cola de un resultado tan grande o mayor bajo $H_{0}$.
\end{frame}

\begin{frame}[fragile]{Paso 2: Valor $p$ y decisión}
  \begin{columns}[T]
    \begin{column}{0.40\textwidth}
      \small
      \begin{lstlisting}
p = stats.norm.sf(z)
reject = p < 0.05
      \end{lstlisting}

      \begin{block}{Salida verificada}
        $p=0.009382$\\
        $\alpha=0.05$\\
        Rechazar $H_{0}$
      \end{block}
    \end{column}
    \begin{column}{0.60\textwidth}
      \centering
      \includegraphics[width=\textwidth]{../figuras/prueba_z_media_02_valor_p.png}
      \vspace{0.2em}
      \scriptsize Cola de la normal estándar del Paso 2
    \end{column}
  \end{columns}
\end{frame}

\begin{frame}{Paso 3: $\sigma$ conocida rara vez es verdadera}
  \begin{columns}[T]
    \begin{column}{0.42\textwidth}
      \small
      $s$ muestral $=4.955099\neq 6$.

      \vspace{0.4em}
      z inválido con $s$:
      \[
        z=2.845789,\quad p=0.002215.
      \]

      \vspace{0.5em}
      \textbf{Límite:} conservar una curva $z$ después de reemplazar $\sigma$
      por $s$ es la prueba incorrecta. La Lección 35 usa $t$ con
      $\mathrm{gl}=n-1$.
    \end{column}
    \begin{column}{0.58\textwidth}
      \centering
      \includegraphics[width=\textwidth]{../figuras/prueba_z_media_03_limite_sigma_conocida.png}
      \vspace{0.2em}
      \scriptsize z válido frente a z inválido con $s$ del Paso 3
    \end{column}
  \end{columns}
\end{frame}

\begin{frame}{Verificación de conceptos}
  \begin{enumerate}
    \item Calcule $\mathrm{EE}=6/\sqrt{40}$.
    \item Calcule $z=(82.229590-80)/\mathrm{EE}$.
    \item ¿Por qué $p=0.009382$ no es la probabilidad de que $\mu=80$?
    \item ¿Por qué $s=4.955099$ no autoriza a conservar la prueba z?
  \end{enumerate}

  \vspace{0.6em}
  \begin{block}{Conexión con la Lección 35}
    La sesión puede continuar directamente con la prueba $t$ de una muestra,
    con $s$ en el denominador y $\mathrm{gl}=n-1$.
  \end{block}

  \vfill
  \scriptsize Fuente: OpenStax, \textit{Introductory Business Statistics 2e},
  Capítulo 9, Secciones 9.3--9.4. CC BY-NC-SA.
\end{frame}

\appendix



\end{document}
